\chapter{Construcción de Paley, estructura compleja finita y código ternario
extendido de Golay}
\label{chap:paley-codigo-witt}
\index{Paley!matriz de conferencia}
\index{Golay!código ternario extendido}
\index{Witt!sistema de Steiner}

La prolongación excepcional de la frontera nonádica requiere una construcción
algebraica anterior a cualquier uso de hexadas, octadas o grupos de Mathieu.
Esta sección fija esa construcción en coordenadas. El orden lógico es
\[
 (\mathbb Z/9\mathbb Z)^\times
 \longrightarrow \mathbb P^1(\mathbb F_5)
 \longrightarrow C_P
 \longrightarrow A_W
 \longrightarrow C_W
 \longrightarrow W_{12}.
\]
Las flechas no se deducen de coincidencias numéricas: las dos primeras fijan
un sistema de coordenadas, las dos siguientes son operaciones matriciales y la
última toma soportes proyectivos de peso seis.

\section{Coordenadas proyectivas y matriz de conferencia}

El conjunto de unidades módulo nueve tiene seis elementos. Elegimos la
bijección de coordenadas
\[
\begin{array}{c|cccccc}
u&1&2&4&8&7&5\\ \hline
\theta(u)&\infty&0&1&2&3&4
\end{array}
\]
en la que la primera fila sigue las potencias del generador \(2\) módulo
nueve. Esta elección fija un origen y una orientación, e identifica el
conjunto subyacente con la recta proyectiva
\[
 \mathbb P^1(\mathbb F_5)=\{\infty,0,1,2,3,4\}.
\]
La aplicación \(\theta\) es un calibre de coordenadas, no un homomorfismo ni
una consecuencia canónica de la estructura multiplicativa de las unidades.
El carácter cuadrático de \(\mathbb F_5\) determina, en el orden anterior, la
matriz simétrica de conferencia
\begin{equation}
 C_P=
 \begin{pmatrix}
 0& 1& 1& 1& 1& 1\\
 1& 0& 1&-1&-1& 1\\
 1& 1& 0& 1&-1&-1\\
 1&-1& 1& 0& 1&-1\\
 1&-1&-1& 1& 0& 1\\
 1& 1&-1&-1& 1& 0
 \end{pmatrix}.
 \label{eq:paley-conference-explicita}
\end{equation}
El producto fila por fila da
\begin{equation}
 C_PC_P^{\mathsf T}=5I_6.
 \label{eq:paley-conference-identidad}
\end{equation}
La diagonal nula y la ortogonalidad de filas expresan, respectivamente, la
ausencia de autoincidencia y el equilibrio del carácter cuadrático. Un cambio
simultáneo de orden o de signos produce una matriz monomialmente equivalente;
la matriz \eqref{eq:paley-conference-explicita} fija el calibre empleado en
todo el manuscrito.

\section[Reducción ternaria y estructura compleja interna]{Reducción
ternaria y endomorfismo cuadrático
\texorpdfstring{\(A_W^2=-I_6\)}{AW2=-I6}}

Reducimos las entradas de \(C_P\) módulo tres, identificando
\(-1\equiv2\pmod3\). Se obtiene
\begin{equation}
 A_W=
 \begin{pmatrix}
 0&1&1&1&1&1\\
 1&0&1&2&2&1\\
 1&1&0&1&2&2\\
 1&2&1&0&1&2\\
 1&2&2&1&0&1\\
 1&1&2&2&1&0
 \end{pmatrix}
 \in M_6(\mathbb F_3).
 \label{eq:AW-explicita}
\end{equation}
Como \(A_W=A_W^{\mathsf T}\), la reducción de
\eqref{eq:paley-conference-identidad} implica
\begin{equation}
 A_W^2=A_WA_W^{\mathsf T}=2I_6=-I_6.
 \label{eq:AW-raiz-menos-uno}
\end{equation}
Por tanto \(A_W\) realiza una raíz cuadrada de \(-I\) dentro de un espacio
vectorial definido sobre \(\mathbb F_3\). No se identifica el cuerpo ternario
con \(\mathbb C\); se afirma la identidad matricial exacta
\eqref{eq:AW-raiz-menos-uno}, que proporciona una estructura compleja finita
en el sentido algebraico de un endomorfismo \(J\) con \(J^2=-I\).

\begin{proposition}[Descomposición en subespacios propios conjugados]
Sea \(V=\mathbb F_3^6\) y sea \(J=A_W\). La extensión escalar
\(V\otimes_{\mathbb F_3}\mathbb F_9\) se descompone en los dos espacios
propios de \(J\) correspondientes a las raíces \(\iota\) y \(-\iota\) de
\(X^2+1\) en \(\mathbb F_9\).
\end{proposition}

\begin{proof}
El polinomio \(X^2+1\) no tiene raíces en \(\mathbb F_3\), pues los
cuadrados son cero y uno. En \(\mathbb F_9\) se factoriza como
\((X-\iota)(X+\iota)\), con raíces distintas. La identidad \(J^2=-I\)
muestra que el polinomio mínimo de \(J\) divide a \(X^2+1\); como \(J\) no
es escalar, coincide con él. El endomorfismo es, por tanto, diagonalizable
después de extender escala a \(\mathbb F_9\), y la suma de ambos espacios
propios es directa y exhaustiva.
\end{proof}

\section{Código gráfico autodual}

Definimos el morfismo lineal
\[
 c:\mathbb F_3^6\longrightarrow\mathbb F_3^{12},
 \qquad c(q)=(q,qA_W),
\]
y su imagen
\[
 C_W=\operatorname{im}c.
\]
La matriz generadora es \(G=(I_6\mid A_W)\). Por
\eqref{eq:AW-raiz-menos-uno},
\[
 GG^{\mathsf T}=I_6+A_WA_W^{\mathsf T}=I_6+2I_6=0
 \quad\text{en }\mathbb F_3.
\]
Las seis filas de \(G\) son independientes, de modo que \(C_W\) tiene
dimensión seis y es igual a su dual ortogonal. La enumeración de los
\(3^6=729\) vectores \(q\) proporciona el enumerador de pesos
\begin{equation}
 W_{C_W}(y)=1+264y^6+440y^9+24y^{12}.
 \label{eq:enumerador-CW}
\end{equation}
En particular, la distancia mínima es seis y
\[
 C_W=[12,6,6]_3,
\]
el código ternario extendido de Golay en las coordenadas fijadas por
\eqref{eq:AW-explicita}.

La enumeración exhaustiva de los \(729\) mensajes, formando cada palabra con
la matriz impresa y contando sus pesos, produce
\eqref{eq:enumerador-CW}.

\begin{proposition}[Palabra marcada y controles uniformes]
\label{prop:palabra-marcada-witt}
En la presentación orientada determinada por el corte \(6+6\), la matriz
\(A_W\) de \eqref{eq:AW-explicita} y la orientación positiva, sea
\[
q_\ast=(1,1,1,1,1,1)\in\mathbb F_3^6.
\]
Entonces
\[
\boxed{
c_\ast=c(q_\ast)
=(1,1,1,1,1,1,2,1,1,1,1,1)\in C_W.}
\]
En cambio,
\[
(1,\ldots,1)\notin C_W,
\qquad
(2,\ldots,2)\notin C_W.
\]
La primera palabra es canónica relativamente a la presentación orientada
declarada; no es una palabra absoluta independiente del calibre.
\end{proposition}

\begin{proof}
La suma de las filas de la matriz impresa da
\[
q_\ast A_W=(2,1,1,1,1,1),
\]
que prueba la fórmula de \(c_\ast\). Para las dos palabras uniformes, la
primera mitad fija respectivamente \(q=q_\ast\) y \(q=-q_\ast\), mientras
las segundas mitades requeridas serían \(q_\ast\) y \(-q_\ast\); las
segundas mitades efectivas son
\((2,1,1,1,1,1)\) y \((1,2,2,2,2,2)\). Por tanto ninguna palabra uniforme
pertenece al código. La enumeración exhaustiva de las \(729\) palabras
confirma que no aparece otra excepción.
\end{proof}

\[
\mathsf R=
\begin{pmatrix}
0&1\\
1&0
\end{pmatrix},
\qquad
\mathsf R C\mathsf R=C^{-1}.
\]
La palabra marcada especializa, después de la construcción abstracta del
\cref{thm:funtor-ciclotomico-fase-longitud}, la familia de trivializaciones
mediante
\[
c_b:=c_{b+1},
\qquad c\in C_W,\quad b\in\{0,\ldots,11\},
\]
donde la expresión de la derecha usa la numeración impresa
\(1,\ldots,12\). Esta convención se mantiene en todo lo que sigue y tipa, en
particular, \((c_\ast)_b\) y el posterior \(P_b(c)\). Entonces
\[
P_b^{(\ast)}=
\begin{cases}
\mathsf R,&(c_\ast)_b=1,\\
I,&(c_\ast)_b=2.
\end{cases}
\]
Esta regla depende del corte, de la matriz y de la orientación declarados;
no interviene en la definición previa del funtor fase--longitud.

La ausencia de las palabras uniformes es un control de esta orientación:
impide usarlas en la prueba concreta del vecino, pero no altera por sí sola
el pegado Niemeier, que se conserva por la acción trivial de cada Coxeter
sobre \(A_2^\ast/A_2\).

\section{Soportes proyectivos y sistema de Witt}

Cada una de las 264 palabras de peso seis tiene un soporte de seis
coordenadas. Las palabras opuestas \(u\) y \(-u\) poseen el mismo soporte y
no hay otra identificación, por lo que se obtienen
\[
 \frac{264}{2}=132
\]
hexadas distintas. Denotemos por \(\mathcal H_W\) la familia de esos
soportes.

\begin{theorem}[Sistema de Witt ternario]
El sistema de incidencia
\[
 \bigl(\{1,\ldots,12\},\mathcal H_W\bigr)
\]
es un sistema de Steiner \(S(5,6,12)\). En consecuencia, es una realización
del diseño de Witt \(W_{12}\).
\end{theorem}

\begin{proof}
El recorrido de los \(\binom{12}{5}=792\) subconjuntos de cinco puntos cuenta
las hexadas que contienen cada uno. Todos los conteos valen
uno. Ésta es exactamente la propiedad definitoria de un
\(S(5,6,12)\). El número de bloques también resulta de la identidad de
incidencias
\[
 |\mathcal H_W|\binom65=\binom{12}{5},
\]
que da \(|\mathcal H_W|=132\).
\end{proof}

Esta sección fija por tanto, antes de utilizarla, la cadena completa
\[
 C_P\longmapsto A_W\longmapsto C_W\longmapsto W_{12}.
\]
Las secciones posteriores añaden el marco dodecafásico, las involuciones de
estrella, la elevación relativa a una dodecada binaria y la construcción
reticular. Ninguno de esos pasos modifica el calibre matricial ya establecido.
